% Figure 2: original environment records, with fonts matching Figure 8(a).
% Run code/export_figure2_data.py, then compile this file with XeLaTeX.
% Copy the resulting PDF to 图3.pdf, which the paper already includes.
\documentclass[tikz,border=0pt]{standalone}
\usepackage{fontspec}
\usepackage{xeCJK}
\usepackage{pgfplots}
\setmainfont{Times New Roman}
\setCJKmainfont{SimSun}
\xeCJKDeclareCharClass{CJK}{"2103}
\pgfplotsset{compat=1.18}
\definecolor{TemperatureRed}{HTML}{D55E00}
\definecolor{ConcentrationBlue}{HTML}{0072B2}
\definecolor{ReferenceGray}{HTML}{666666}
\def\TemperatureMinimum{26.887699999999999}
\def\TemperatureMaximum{51.3583}
\def\ConcentrationMinimum{0.018099000000000001}
\def\ConcentrationMaximum{0.051781000000000001}

\pgfplotsset{
  environment panel/.style={
    scale only axis,
    width=357.12bp,height=115.92bp,
    xmin=0,xmax=4.8,
    xtick={0,1,2,3,4},xticklabels={0,1,2,3,4},
    axis lines=box,
    axis line style={black,line width=0.8bp},
    tick align=outside,
    xtick pos=bottom,ytick pos=left,
    tick style={black,line width=0.8bp},
    major tick length=3.5bp,
    tick label style={font=\fontsize{10.5}{12}\selectfont},
    label style={font=\fontsize{10.5}{12}\selectfont},
    grid=major,grid style={draw=black!15,line width=0.5bp},
    axis background/.style={fill=white},
    scaled ticks=false,
    clip=true,
  },
}
\begin{document}
\begin{tikzpicture}[x=1bp,y=1bp]
  % Preserve the former PDF's canvas and both axis rectangles.
  \path[use as bounding box] (0,0) rectangle (413.9575,299.549);
  \begin{pgfinterruptboundingbox}
    \begin{axis}[
      environment panel,
      at={(49.6375bp,176.429bp)},anchor=south west,
      ymin=\TemperatureMinimum,ymax=\TemperatureMaximum,
      ytick={30,35,40,45,50},yticklabels={30,35,40,45,50},
      xticklabels={},
      ylabel={烘房温度(℃)},
      ylabel style={at={(axis description cs:-0.077,0.5)},anchor=south},
      legend style={
        at={(0.997,0.060)},anchor=south east,
        fill=white,fill opacity=1,text opacity=1,draw=none,
        font=\fontsize{7}{8.2}\selectfont,
        inner sep=0.8bp,row sep=0.5bp,column sep=1bp,
      },
      legend image code/.code={
        \draw[#1] plot coordinates {(0bp,0bp) (6bp,0bp)};
      },
      legend cell align=left,
    ]
      \addplot[TemperatureRed,line width=1.4bp]
        table[x=time,y=temperature] {fig01_env_records.dat};
      \addlegendentry{附件1实测}
      \addplot[TemperatureRed,line width=1.4bp,
        dash pattern=on 5.18bp off 2.24bp,
        legend image post style={dash pattern=on 2bp off 1.5bp}]
        coordinates {(4,50) (4.8,50)};
      \addlegendentry{恒温段设计值}
      \addplot[ReferenceGray,line width=0.9bp,
        dash pattern=on 0.9bp off 1.485bp,forget plot]
        coordinates {(4,\TemperatureMinimum) (4,\TemperatureMaximum)};
      \node[anchor=base,inner sep=0pt,text=ReferenceGray,
        font=\fontsize{9.5}{11}\selectfont]
        at (axis cs:1.9,30.5) {预热平衡阶段};
      \node[anchor=base,inner sep=0pt,text=ReferenceGray,
        font=\fontsize{9.5}{11}\selectfont]
        at (axis cs:4.4,44.5) {恒温干燥阶段};
    \end{axis}
    \begin{axis}[
      environment panel,
      at={(49.6375bp,37.325bp)},anchor=south west,
      ymin=\ConcentrationMinimum,ymax=\ConcentrationMaximum,
      ytick={0.02,0.03,0.04,0.05},
      yticklabels={0.02,0.03,0.04,0.05},
      ylabel={烘房水分浓度(kg/kg)},
      ylabel style={at={(axis description cs:-0.077,0.5)},anchor=south},
      xlabel={时间(h)},
      xlabel style={at={(axis description cs:0.5,-0.16)},anchor=north},
    ]
      \addplot[ConcentrationBlue,line width=1.4bp]
        table[x=time,y=concentration] {fig01_env_records.dat};
      \addplot[ConcentrationBlue,line width=1.4bp,
        dash pattern=on 5.18bp off 2.24bp]
        coordinates {(4,0.05) (4.8,0.05)};
      \addplot[ReferenceGray,line width=0.9bp,
        dash pattern=on 0.9bp off 1.485bp]
        coordinates {(4,\ConcentrationMinimum) (4,\ConcentrationMaximum)};
    \end{axis}
  \end{pgfinterruptboundingbox}
\end{tikzpicture}
\end{document}
